\documentclass{article}
% Stable PDF metadata; the visible review copy remains anonymous.
\pdfinfoomitdate=1
\pdftrailerid{}
\PassOptionsToPackage{numbers,sort&compress}{natbib}
\usepackage[dblblindworkshop,nonanonymous]{neurips_2026}
\usepackage[T1]{fontenc}
\usepackage[utf8]{inputenc}
\usepackage{amsmath,amssymb,amsthm,mathtools}
\usepackage{booktabs,array,multirow}
\usepackage{microtype}
\usepackage{graphicx}
\usepackage{tikz,pgfplots}
\pgfplotsset{compat=1.18}
\usetikzlibrary{arrows.meta,positioning,fit,calc}
\usepackage{url}
\usepackage[hidelinks]{hyperref}
\usepackage{aliascnt}
\usepackage[nameinlink,noabbrev]{cleveref}
\newtheorem{theorem}{Theorem}
\newaliascnt{proposition}{theorem}
\newtheorem{proposition}[proposition]{Proposition}
\aliascntresetthe{proposition}
\crefname{proposition}{proposition}{propositions}
\Crefname{proposition}{Proposition}{Propositions}
\newaliascnt{corollary}{theorem}
\newtheorem{corollary}[corollary]{Corollary}
\aliascntresetthe{corollary}
\crefname{corollary}{corollary}{corollaries}
\Crefname{corollary}{Corollary}{Corollaries}
\newaliascnt{lemma}{theorem}
\newtheorem{lemma}[lemma]{Lemma}
\aliascntresetthe{lemma}
\crefname{lemma}{lemma}{lemmas}
\Crefname{lemma}{Lemma}{Lemmas}
\theoremstyle{definition}
\newaliascnt{definition}{theorem}
\newtheorem{definition}[definition]{Definition}
\aliascntresetthe{definition}
\crefname{definition}{definition}{definitions}
\Crefname{definition}{Definition}{Definitions}
\newaliascnt{example}{theorem}
\newtheorem{example}[example]{Example}
\aliascntresetthe{example}
\crefname{example}{example}{examples}
\Crefname{example}{Example}{Examples}
\newcommand{\C}{\mathcal C}
\newcommand{\Latt}{\mathcal L}
\newcommand{\R}{\mathcal R}
\newcommand{\V}{\mathcal V}
\newcommand{\NN}{\mathbb N}
\newcommand{\RR}{\mathbb R}
\newcommand{\ind}{\mathbf 1}
\newcommand{\relu}{\operatorname{ReLU}}
\newcommand{\height}{\operatorname{ht}}
\newcommand{\Idem}{\operatorname{Idem}}
\newcommand{\Con}{\operatorname{Con}_{\wedge}}
\newcommand{\Mcomp}{M_{\mathrm{comp}}}
\newcommand{\Mstat}{M_{\mathrm{stat}}}
\newcommand{\dec}[1]{[\![#1]\!]}
\newcommand{\argmax}{\operatorname*{arg\,max}}
\newcommand{\argmin}{\operatorname*{arg\,min}}
\newcommand{\eps}{\varepsilon}
\newcommand{\powerset}{\mathcal P}
\newcommand{\code}{\mathsf{code}}
\newcommand{\compile}{\mathsf{compile}}
\newcommand{\extract}{\mathsf{extract}}
\newcommand{\cl}{\operatorname{cl}}
\newcommand{\kerop}{\operatorname{ker}}
\newcommand{\Min}{\operatorname{Min}}
\newcommand{\sym}{\mathsf{sym}}
\newcommand{\obs}{\mathsf{obs}}
\newcommand{\model}{\mathsf{model}}
\newcommand{\doop}{\operatorname{do}}
\newcommand{\readout}{\rho}
\newcommand{\bottomans}{\mathsf{inconsistent}}
\newcommand{\unknown}{\mathsf{undetermined}}
\newcommand{\ceil}[1]{\left\lceil #1\right\rceil}
\newcommand{\norm}[1]{\left\lVert #1\right\rVert}
\newcommand{\abs}[1]{\left|#1\right|}
\hypersetup{pdftitle={Model Inference from Neural Artifacts: An Exchange Law for Structure and Memory},pdfauthor={},pdfsubject={Neural artifact analysis and finite-state model inference},pdfkeywords={neural artifacts, model inference, version spaces, attention, model extraction}}
\title{Model Inference from Neural Artifacts:\\An Exchange Law for Structure and Memory}
\author{Anonymous Author(s)}
\workshoptitle{Neural Network Artifacts as a New Data Modality}
\begin{document}
\maketitle
\begin{abstract}
An artifact's inference state has two capacities, cardinality and chain height, and independent pooling is priced by the second. The state is the set of distinctions that future observations can expose, a finite meet-semilattice $\R$. Coupled attention needs $\lceil\log_b|\R|\rceil$ coordinates of $b$ values; independently pooled channels of sizes $a_1,\ldots,a_d$ need $\sum_j(a_j-1)\ge\height(\R)-1$, for arbitrary commutative pooling and redundant implementation states. On an $n$-state chain the two are exact and separate exponentially, $\lceil(n-1)/(b-1)\rceil$ against $\lceil\log_b n\rceil$. On a Boolean store of $k$ independently eliminated candidates they coincide at $k$. Chain height is the number of times future evidence can change the answer, so on a pooled memory adaptive queries and memory are one resource: a policy whose memory is the store makes at most $\sum_j(a_j-1)$ informative queries, whereas a decoder with unrestricted memory exchanges stored bits and query bits one for one. Attention--ReLU artifacts realize both sides of each law, their minimal inference machines are recovered from operational parameters, and finite-symbol hardening preserves the machine whenever the probability mass on wrong symbols lies below an explicit threshold. Most results are machine-checked in Lean~4.\footnote{Lean sources, paper source, and checks: \url{https://anonymous.4open.science/r/neural-artifacts-exchange-law-838F/}}
\end{abstract}

\section{Introduction}
Neural artifacts contain more than the answers returned on an evaluation set. Their weights, routing decisions, and intermediate representations also determine how information is stored, combined, and used after new evidence arrives. Learning directly on network parameters already supports model comparison and prediction, with parameter symmetries providing a central architectural principle \citep{zhou2023nfn}. We study a complementary question: \emph{which model of inference can be recovered from an artifact, and which distinctions must its internal representations preserve?}

A network can represent candidate models through different mechanisms. Attention can expose a selection rule or a compositional routing tree; feedforward computations can retain tables, transformations, and residual information. These roles do not assign an intrinsic meaning to a layer type. Instead, we specify their common semantic interface: each representation participates in maintaining the model possibilities consistent with observations. The central issue is then whether information stored in one representation can replace information stored in another while preserving future inference.

Counting states gives a necessary condition, but misses the cost of composition. A four-state chain can be encoded by two bits. It cannot be maintained by two independently pooled binary channels, even when their pooling operations are chosen freely. Three such channels suffice. The discrepancy is not a shortage of bits at a single instant: it arises because each channel must combine evidence without consulting the others. We prove the exact law for every chain size and every finite commutative pooling architecture. At fixed state cardinality, a Boolean store has a different exchange law because complementary coordinates preserve intersection independently.

Our semantic starting point is a version space \citep{mitchell1982generalization}. An observation removes incompatible candidates, and an inference rule reads the survivors. We identify two version spaces when no permitted continuation changes the inference rule's answer between them. This is the classical automata-minimization construction \citep{moore1956,nerode1958}, applied to the meet-semilattice generated by model observations. Every correct implementation has a unique interpretation onto this quotient on its reachable states. The inference task forces the semantic object.

The inference state $\R$ forced by a task carries two numbers, $|\R|$ and $\height(\R)$. Coupled pooling is priced by the first and independent pooling by the second, so the exchange rate between structure and memory is a property of the store, and the separation of \cref{cor:width} appears exactly when $\height(\R)$ approaches $|\R|$. Chain height is mind-change complexity, the number of times a readout can change along a history; a store that affords many revisions at logarithmic width must couple its coordinates. The same quantity bounds adaptive experimentation on a pooled memory (\cref{thm:policy}), which places the memory law of \cref{app:tables} and the exchange law on one scale. Finite-state extraction and an explicit attention--ReLU compiler connect the laws to artifacts, and finite-symbol hardening preserves the extracted machine, including ties between branches with the same symbol.

\paragraph{Position relative to prior work.}
Knowledge compilation compares representations through their succinctness and supported queries and transformations \citep{darwiche2002map}. We apply this perspective to the persistent state of neural inference and derive sharp costs for independently compositional representations. The kernel criterion for exact factorization is the classical subdirect-product criterion \citep{birkhoff1944}; the new quantitative result is the exact exchange region for pooled inference, together with matching neural constructions. Automata extraction from recurrent networks uses queries and abstractions to recover state dynamics \citep{weiss2018extracting}. Our extraction theorem supplies an exact bridge from a finite neural execution interface to the same canonical object used by the lower bounds. Soft attention can simulate hard attention under score-gap conditions and preserves a final classifier \citep{yang2026soft}; finite-symbol correction preserves the entire model-inference machine. Slot Attention provides an important example of exchanging information between perceptual features and object-level representations \citep{locatello2020slot}. The algebraic interface below specifies what must be verified to give such an exchange an exact model-inference semantics. Chain height is the mind-change complexity of inductive inference \citep{gold1967,casesmith1983}; the exchange law prices it by the pooling operation.

\section{The inference state represented by an artifact}\label{sec:inference}
Let $\C$ be a finite set of candidate models, $P$ an observation alphabet, and $K:P\to\powerset(\C)$ a consistency map. The alphabet can include labeled examples, experimental results, or inspected artifact properties. For a history $s=p_1\cdots p_t$, define
\begin{equation}
 V(s)=\bigcap_{i=1}^tK(p_i),\qquad V(\epsilon)=\C,
 \qquad \Latt=\{V(s):s\in P^*\}.
 \label{eq:version}
\end{equation}
The finite family $\Latt$ contains $\C$ and is closed under intersection, since $V(st)=V(s)\cap V(t)$. Fix a readout $\rho:\Latt\to Y$. First-consistent-model selection and query certification are examples. Empty version spaces receive an explicit inconsistency answer whenever consistency is relevant.

\begin{definition}[Continuation equivalence]\label{def:nerode}
For $U,V\in\Latt$, let
\begin{equation}
 U\equiv_\rho V
 \quad\Longleftrightarrow\quad
 \rho(U\cap H)=\rho(V\cap H)\quad\text{for every }H\in\Latt.
 \label{eq:nerode}
\end{equation}
Write $\R=\Latt/{\equiv_\rho}$, with $[U]\wedge[V]=[U\cap V]$ and top element $[\C]$.
\end{definition}
The relation respects intersection: $H\cap W\in\Latt$, so equivalence of $U,V$ implies equivalence of $U\cap W,V\cap W$. Taking $H=\C$ shows that the readout descends to $\R$.

\begin{theorem}[Canonical inference state]\label{thm:canonical}
The least number of states of a deterministic machine computing $s\mapsto\rho(V(s))$ is $|\R|$. The same minimum holds for machines that pool observations in a finite commutative monoid. Every correct deterministic implementation has a unique map $\delta$ from its reachable states onto $\R$ satisfying
\begin{equation}
 \delta(m_\epsilon)=[\C],\qquad
 \delta(\operatorname{step}(m,p))=\delta(m)\wedge[K(p)],
 \qquad \operatorname{out}(m)=\bar\rho(\delta(m)).
 \label{eq:simulation}
\end{equation}
\end{theorem}
\begin{proof}
For a state reached by $s$, set $\delta(m_s)=[V(s)]$. If $m_s=m_t$, every common continuation produces the same output. Every $H\in\Latt$ is the version space of a continuation, so $V(s)\equiv_\rho V(t)$. This proves well-definedness. Every quotient class is reachable; the update and readout equations follow from intersection. Induction on histories proves uniqueness. Hence any implementation has at least $|\R|$ reachable states. Conversely, initialize at $[\C]$, pool $[K(p)]$ by meet, and apply $\bar\rho$. This commutative machine has exactly $|\R|$ states.
\end{proof}

The quotient depends on both the available observations and the required answers. A constant readout gives one state even when $\Latt$ is large. An injective readout gives $\R=\Latt$. Enlarging the observation language or enriching the readout changes the distinctions that an artifact must retain. In particular, a width-$W$ state with $b$ possible values per coordinate satisfies
\begin{equation}
 |\R|\le b^W.\label{eq:cardinality}
\end{equation}
Here $b$ is an alphabet size: $p$-bit coordinates have $b=2^p$. Equation~\eqref{eq:cardinality} counts persistent inference states, independently of the implementation of the candidate models.

The second number is the height $\height(\R)$, the number of elements in a longest chain. Along any history the states $[V(s)]$ descend in $\R$, so the readout changes at most $\height(\R)-1$ times. Cardinality is what a coupled store must hold; height is what an independent one must hold (\cref{sec:exchange}).

\paragraph{Concrete stores.}
Nested consistency sets $U_0\subsetneq\cdots\subsetneq U_{n-1}=\C$, together with a readout that distinguishes them, give the chain $C_n$. Positive observations of ordered thresholds give this form with a suitable domain and readout. Independent candidate eliminations generate the Boolean store $2^{[k]}$; reading the full survivor set distinguishes all $2^k$ states.

\section{An exact exchange law}\label{sec:exchange}
An independent pooled channel consists of a finite commutative monoid $M_j$ and an observation embedding $e_j:P\to M_j$. Its state after $s$ is
\begin{equation}
 \mu_j(s)=\prod_{p\in s}e_j(p).
 \label{eq:independent}
\end{equation}
A joint readout can inspect every channel, but each channel combines observations using its own operation. The monoids need not be idempotent, and different histories with the same inference state can have different implementations. We write $\height(Q)$ for the number of elements in a longest chain of a finite partially ordered set $Q$.

\begin{theorem}[Exchange under independent pooling]\label{thm:exchange}
Suppose $d$ independent pooled channels compute the inference task of \cref{sec:inference}. Let $E_j=\Idem(M_j)=\{x:x^2=x\}$, ordered by $x\le y$ when $xy=x$. Then
\begin{equation}
 |\R|\le\prod_{j=1}^d|M_j|,
 \qquad
 \height(\R)-1\le\sum_{j=1}^d\bigl(\height(E_j)-1\bigr)
                 \le\sum_{j=1}^d(|M_j|-1).
 \label{eq:height}
\end{equation}
For $\R=C_n$ and channel budgets $a_j\ge1$, a correct implementation with $|M_j|\le a_j$ exists if and only if
\begin{equation}
 \boxed{\ \sum_{j=1}^d(a_j-1)\ge n-1.\ }\label{eq:chainregion}
\end{equation}
\end{theorem}
\begin{proof}
Let $S$ be the image of $s\mapsto(\mu_1(s),\ldots,\mu_d(s))$. It is a submonoid of $\prod_jM_j$. By \cref{thm:canonical}, the map $\delta:S\to\R$ is well-defined and surjective. It is a monoid homomorphism because concatenation corresponds to multiplication in $S$ and meet in $\R$.

Every element of a finite monoid has an idempotent positive power. Since $\delta(x^r)=\delta(x)$ for $r\ge1$, the restriction $\delta:\Idem(S)\to\R$ remains surjective. Commutativity makes $\Idem(S)$ a meet-semilattice, contained in $\prod_jE_j$. Lift a chain $r_1<\cdots<r_h$ in $\R$ to idempotents $e_i$ and set $z_i=e_i\cdots e_h$. Then $z_i\le z_{i+1}$ and $\delta(z_i)=r_i$, so the lifted chain is strict. A chain in a product has at most $1+\sum_j(\height(E_j)-1)$ elements: every strict step advances at least one coordinate. This proves \eqref{eq:height} and the necessity of \eqref{eq:chainregion}.

For sufficiency, label the chain $0<\cdots<n-1$. Distribute its $n-1$ adjacent boundaries among sets $B_j$, with $|B_j|\le a_j-1$. The map
\begin{equation}
 f_j(r)=|\{c\in B_j:c<r\}|\label{eq:cuts}
\end{equation}
preserves minimum and takes $|B_j|+1$ values. Every boundary is distinguished by some coordinate, so the paired map is injective. Pool with minimum in each coordinate and decode the joint state.
\end{proof}

The use of idempotent powers is essential: the lower bound applies to redundant encodings, cyclic state components, and all other finite commutative pooling operations. In particular, a finite group contributes no order height because its identity is its only idempotent.

\begin{corollary}[Independent versus coupled width]\label{cor:width}
For an $n$-state chain and $b\ge2$, the minimum number of independently pooled $b$-value coordinates is
\begin{equation}
 W_{\mathrm{ind}}(n,b)=\ceil{\frac{n-1}{b-1}}.
 \label{eq:widthind}
\end{equation}
If the pooling operation can mix coordinates, the exact minimum is
\begin{equation}
 W_{\mathrm{joint}}(n,b)=\ceil{\log_b n}.\label{eq:widthjoint}
\end{equation}
\end{corollary}
\begin{proof}
The first equality is \cref{thm:exchange}. For the second, encode the rank in base $b$ and implement minimum on the represented ranks. The cardinality bound \eqref{eq:cardinality} is necessary and attained. Place unused codes below the reachable chain to extend minimum to the full $b^W$-element carrier. The encoded top remains the monoid identity and unused codes are not reached.
\end{proof}

For $n=2^k$, independent binary pooling requires $2^k-1$ persistent coordinates, while coupled pooling requires $k$. A coupled head can implement minimum directly: assign keys $-r,-s$, query $1$, and values equal to the digit encodings of $r,s$. Hard attention selects the smaller rank. For equal ranks the values coincide, so any tie rule gives the same result.

\begin{corollary}[Width for an arbitrary inference state]\label{cor:general}
For any finite $\R$ and $b\ge2$,
\begin{equation}
 W_{\mathrm{joint}}(\R,b)=\ceil{\log_b|\R|},\qquad
 W_{\mathrm{ind}}(\R,b)\ge\max\left\{\ceil{\log_b|\R|},\ \ceil{\frac{\height(\R)-1}{b-1}}\right\}.
 \label{eq:general}
\end{equation}
For the Boolean store $2^{[k]}$ with survivor readout, $W_{\mathrm{ind}}=W_{\mathrm{joint}}=k$.
\end{corollary}
\begin{proof}
The coupled construction of \cref{cor:width} extends to any $\R$ by sending unused codes to an adjoined absorbing bottom, which keeps the carrier a commutative idempotent monoid with identity $[\C]$; \eqref{eq:cardinality} forbids fewer coordinates. The lower bound is \eqref{eq:cardinality} together with \eqref{eq:height} at $|M_j|\le b$. One binary channel per candidate, pooled by conjunction, reaches every survivor set of $2^{[k]}$ with an injective readout, and $|2^{[k]}|=2^k$ forbids fewer coordinates under any pooling.
\end{proof}
The separation belongs to chains: $W_{\mathrm{ind}}/W_{\mathrm{joint}}$ is controlled by $\height(\R)/\log_b|\R|$.

\subsection{How a structural view replaces residual memory}
A normalized compositional view of $\R$ is a meet homomorphism $a:\R\to A$. Its kernel $\alpha$ is a meet congruence: equivalent states remain equivalent after intersection with any state. A second view $g:\R\to G$, with kernel $\beta$, jointly determines the inference state precisely when
\begin{equation}
 \alpha\cap\beta=\Delta_\R.\label{eq:subdirect}
\end{equation}
This is the subdirect-product criterion \citep{birkhoff1944}. At a fixed structural view, its exact compositional residual requirement is therefore
\begin{equation}
 \Mcomp(\alpha)=
 \min_{\substack{\beta\in\Con(\R)\\\alpha\cap\beta=\Delta_\R}}
 |\R/\beta|.
 \label{eq:comp}
\end{equation}
Refining the structural view from $\alpha$ to $\alpha'\subseteq\alpha$ enlarges the feasible set in \eqref{eq:comp}; hence $\Mcomp(\alpha')\le\Mcomp(\alpha)$. In contrast, an arbitrary residual code used to reconstruct the state once requires exactly
\begin{equation}
 \Mstat(\alpha)=\max_{F\in\R/\alpha}|F|.\label{eq:static}
\end{equation}
Assign distinct labels within each fibre and reuse them across fibres.

\begin{proposition}[Sharp chain and Boolean exchange]\label{prop:exact}
For a chain $C_n$ and any structural meet quotient with $a$ states,
\begin{equation}
 \boxed{\Mcomp(\alpha)=n-a+1.}\label{eq:residualchain}
\end{equation}
The same minimum holds when the residual is an arbitrary finite commutative pooled channel, including redundant histories. For $\R=2^{[k]}$ with structural view $U\mapsto U\cap J$, $|J|=s$,
\begin{equation}
 \boxed{\Mcomp(\alpha)=\Mstat(\alpha)=2^{k-s}.}\label{eq:boolean}
\end{equation}
\end{proposition}
\begin{proof}
A chain congruence has interval blocks: if $x\sim z$ and $x\le y\le z$, meeting with $y$ gives $x\sim y$. Its $a$ blocks distinguish $a-1$ boundaries. The residual must distinguish the other $n-a$ boundaries, and cutting exactly these boundaries attains $n-a+1$ states. The general lower bound follows from \cref{thm:exchange}, with the fixed structural quotient as the first channel. For the Boolean store, every structural fibre has $2^{k-s}$ elements, requiring that many distinct residual values. Projection to $[k]\setminus J$ attains the bound and preserves intersection.
\end{proof}

\begin{figure}[t]
\centering
\begin{tikzpicture}
\begin{axis}[width=.83\linewidth,height=2.9cm,xmin=1,xmax=16,ymin=1,ymax=16,
 xlabel={Structural states $a$},ylabel={Residual states},
 xtick={1,4,8,12,16},ytick={1,4,8,12,16},
 legend style={at={(0.98,0.98)},anchor=north east,font=\small,draw=none},
 tick label style={font=\small},label style={font=\small}]
\addplot[thick,domain=1:16,samples=16]{17-x};
\addlegendentry{Chain: $17-a$}
\addplot[thick,dashed,mark=square*,mark size=2pt] coordinates {(1,16)(2,8)(4,4)(8,2)(16,1)};
\addlegendentry{Boolean coordinate views: $16/a$}
\addplot[only marks,mark=*,mark size=2pt] coordinates {(4,13)};
\node[anchor=south west,font=\small] at (axis cs:4,13) {$(4,13)$};
\node[anchor=north east,font=\small] at (axis cs:4,4) {$(4,4)$};
\end{axis}
\end{tikzpicture}
\caption{Exact structure--memory exchange for two 16-state inference tasks. A four-state structural view leaves 13 compositional residual states in a chain and four in a Boolean store with two exposed coordinates. The Boolean line joins the admissible coordinate projections. State cardinality alone does not determine the exchange rate.}
\label{fig:exchange}
\end{figure}

For the four-state chain, take structural blocks $\{0,1\},\{2,3\}$. A static binary residual can record parity. It cannot update by itself: $0$ and $2$ have the same parity, but meeting each with $1$ gives $0$ and $1$. A compositional residual instead uses $\{0\},\{1,2\},\{3\}$. This example separates recoverable information from independently reusable information.

\section{Recovering models from neural artifacts}\label{sec:neural}
The exchange law concerns the semantic state represented by an artifact. We now give operational extraction and realization procedures. An artifact specifies its architecture, parameters, initial state, numerical semantics, and readout; these determine a transition function $F_\theta(z,p)$.

\begin{theorem}[Finite-state extraction]\label{thm:extract}
Let $Z$ be a finite state codebook containing $z_0$, let $P$ be finite, and suppose $F_\theta:Z\times P\to Z$ and $o_\theta:Z\to Y$ can be evaluated exactly. From the operational artifact one can recover the reachable transition table, its minimal Moore machine, and a state-by-state simulation onto that machine. If the artifact computes $\rho(V(s))$ for every $s\in P^*$, the extracted minimal machine is isomorphic to $\R$. If its minimized input transitions commute and are idempotent, the artifact itself determines its meet-semilattice representation, and any meet law consistent with its transitions is that one.
\end{theorem}
\begin{proof}
Starting at $z_0$, enumerate every transition from each newly reached state. This terminates after at most $|Z||P|$ evaluations. Partition the reachable states by readout, and repeatedly refine by successor-block signatures until stable. Induction on continuation length shows that the final equivalence is agreement on every future output. The quotient and projection give the required machine and simulation. \Cref{thm:canonical} identifies this quotient with $\R$. For the final statement, choose an access word $w_r$ for each minimal state $r$, and define $r\wedge s$ by running $w_s$ from $r$. Commutativity makes this independent of the access word; idempotence makes it a meet. The initial state is its top. \Cref{app:extraction} proves the representation and its uniqueness.
\end{proof}

A finite-precision execution supplies a finite codebook directly, with every persistent register included. The extraction is invariant under function-preserving hidden-neuron permutations because it depends on the transition function, not the ordering of its implementation units. State-coordinate relabelings yield isomorphic extracted machines. These are the semantic symmetries relevant to model inference, complementing parameter-level equivariance \citep{zhou2023nfn}.

The last clause gives a fully operational analysis: extract the transition system, minimize it, test commutativity and idempotence of its input transitions, and recover the meet operation. The candidate realization $\C=\R$, $K(p)=\{r:r\le t_p(\top)\}$ then has survivor sets $\mathord{\downarrow}r$. The artifact determines this operational version-space model; an external model population supplies its application-specific interpretation. For continuous-state artifacts, a finite invariant-region certificate gives the same extraction theorem; rational hard-attention/ReLU maps and semialgebraic regions admit exact certificate checking (\cref{app:extraction}).

\subsection{Attention forests and finite feedforward maps}
We use a typed code grammar with finite input alphabet $X$. A leaf is a finite map $f:X\to Y$. An internal node has children $c_1,\ldots,c_k$ with common finite output alphabet $Y$, affine scores $s_j(x)$ in a declared feature vector $\phi(x)$, and a finite map $g:Y\to Y'$. Its interpretation is
\begin{equation}
 \dec{c}(x)=g\bigl(\dec{c_{j^*(x)}}(x)\bigr),
 \qquad j^*(x)=\min\argmax_j s_j(x).
 \label{eq:code}
\end{equation}
A forest provides multiple such outputs. Fixed attention reads followed by finite tuple maps compose these outputs into a finite-state update or a model-query program. Neither the grammar nor the store requires the candidate models themselves to be linear.

\begin{theorem}[Neural realization and extraction]\label{thm:compiler}
Every finite attention forest of the form \eqref{eq:code} has a masked dot-product attention implementation with shared positionwise ReLU feedforward maps, residual connections, and no normalization. Its operational parameters determine an extracted code satisfying
\begin{equation}
 \dec{\extract(\compile(c))}(x)=\dec{c}(x)
 \quad\text{for every }x\in X.\label{eq:roundtrip}
\end{equation}
Any finite transition system can be realized by a ReLU update, and composing either construction with \cref{thm:extract} recovers its minimal inference machine.
\end{theorem}
\begin{proof}[Construction]
Encode a symbol $y$ as the basis vector $e_y$. A finite map has matrix $H_ge_y=e_{g(y)}$. An affine score $\alpha_j^\top\phi(x)+\beta_j$ is a dot product between query $(\phi(x),1)$ and key $(\alpha_j,\beta_j)$. Hard attention selects the child value, and the feedforward map applies $H_g$. For shared feedforward dispatch, pair a node-identity coordinate $u_i$ and a symbol coordinate $v_j$ using $\relu(u_i+v_j-1)$; this is their conjunction on one-hot inputs. Masks isolate each parent's children. Separate protected, payload, and scratch coordinates preserve features and implement residual overwrites. \Cref{app:compiler} gives the complete layer construction and the inductive proof.

For a transition table $T:[N]\times[m]\to[N]$, the hidden units
\begin{equation}
 h_{ij}(z,p)=\relu(z_i+p_j-1),\qquad
 F(z,p)=\sum_{i,j}e_{T(i,j)}h_{ij}(z,p)\label{eq:ffntable}
\end{equation}
implement $F(e_i,e_j)=e_{T(i,j)}$ exactly. Extraction evaluates these operational maps and reads the routing scores and wiring. No unused source-code identifier is required.
\end{proof}

The compiler realizes both sides of \cref{cor:width}. Independent channels implement the cut maps \eqref{eq:cuts}; a joint attention operation implements minimum on compact digit codes. Compositional structure supports model inference when the extracted maps preserve the relevant operations; a large feedforward table supports the same task by retaining the corresponding transition information.

\section{Hardening preserves inference through finite symbols}\label{sec:hardening}
Changing softmax weights to argmax decisions changes the executed function in general. We instead give an exact preservation condition at the symbolic interfaces of \cref{thm:compiler}. The correction function below is the standard ReLU rounding construction used in soft-attention simulations \citep{yang2026soft}, parameterized by a tolerance $0<\eta<1/2$:
\begin{equation}
 R_\eta(t)=\frac{\relu(t-\eta)-\relu(t-1+\eta)}{1-2\eta}.
 \label{eq:round}
\end{equation}
Apply it coordinatewise before the finite map $H_g$.

\begin{theorem}[Exact symbolic hardening]\label{thm:hardening}
At a gate with scores $s_j$ and one-hot child values $e_{y_j}$, let $j^*$ be the deterministic hard winner and $y^*=y_{j^*}$. For temperature $T>0$, write
\begin{equation}
 p_j(T)=\frac{e^{s_j/T}}{\sum_i e^{s_i/T}},\qquad
 \eps_{\ne y^*}(T)=\sum_{j:y_j\ne y^*}p_j(T).
 \label{eq:wrongmass}
\end{equation}
Then
\begin{equation}
 R_\eta\left(\sum_jp_j(T)e_{y_j}\right)=e_{y^*}
 \quad\Longleftrightarrow\quad \eps_{\ne y^*}(T)\le\eta.
 \label{eq:exactmass}
\end{equation}
Suppose every maximum-score branch returns $y^*$, there are $m_*$ such branches, and every branch with a different symbol has score at most $s_{j^*}-\delta$ for some $\delta>0$. If $k_-$ branches have a different symbol, then
\begin{equation}
 \eps_{\ne y^*}(T)\le\frac{r(T)}{1+r(T)},\qquad
 r(T)=\frac{k_-}{m_*}e^{-\delta/T}.
 \label{eq:margin}
\end{equation}
A forest whose every gate satisfies \eqref{eq:exactmass} on its symbolic inputs computes exactly the same function at temperature $T$ and at hard attention. When used as a state update, it preserves every history and hence the extracted inference quotient and all exchange laws above.
\end{theorem}
\begin{proof}
The coordinate of $y^*$ in the mixture is $1-\eps_{\ne y^*}$, and the other coordinates sum to $\eps_{\ne y^*}$. The map $R_\eta$ is zero on $[0,\eta]$, one on $[1-\eta,1]$, and strictly between them otherwise. This proves \eqref{eq:exactmass}. The total unnormalized weight of different symbols is at most $k_-e^{(s_{j^*}-\delta)/T}$, while the weight of $y^*$ is at least $m_*e^{s_{j^*}/T}$. Their ratio proves \eqref{eq:margin}. Induction over the forest restores an exact symbol after each gate, so errors do not accumulate with depth. Induction over state updates then proves equality on all histories.
\end{proof}

For a unique winner with gap $\delta$ from every other branch, a sufficient condition is
\begin{equation}
 T\le\frac{\delta}{\log((k-1)(1-\eta)/\eta)}\qquad(k\ge2).
 \label{eq:temperature}
\end{equation}
The symbol-based criterion is stronger than unique-route stability: tied or changing routes cause no loss when they return the same symbol. Conversely, branches with different values at an unresolved tie can invalidate hardening. Without correction, a soft gate mixing constants $0,1$ returns an interior value that a hard selection of those experts does not return.

\paragraph{Execution and identification are different.}
Let $O_T(\theta)$ expose routing probabilities and let $O_0(\theta)=h(O_T(\theta))$ expose only their winners, with other observed quantities held fixed. The observer's version spaces satisfy
\begin{equation}
 \{\theta':O_T(\theta')=O_T(\theta)\}
 \subseteq
 \{\theta':O_0(\theta')=O_0(\theta)\}.\label{eq:coarsen}
\end{equation}
Thus hardening can preserve execution while a coarsened trace reveals less about the artifact. More generally, a target property $Q(\theta)$ can be inferred exactly from behavioral observations $O(\theta)$ and an artifact statistic $A(\theta)$ if and only if
\begin{equation}
 \kerop(O,A)\subseteq\kerop Q.\label{eq:sufficiency}
\end{equation}
Artifact access improves identification precisely when it separates a target-relevant observational equivalence class. Which equivalence class is separated depends on the observation channel.

\section{Model queries, memory, and interventions}\label{sec:queries}
A query $q:\C\to A$ can be certified before a unique model is identified. For nonempty $V$, return $a$ when every $c\in V$ has $q(c)=a$, and $\unknown$ otherwise. Return $\bottomans$ on $V=\varnothing$. This defines another readout in \cref{sec:inference} and therefore another canonical quotient.

\begin{proposition}[Exact query certificate]\label{prop:cover}
Fix a true model $c^*\in V$ and deterministic experiments with outcomes $D(c,e)$. Let
\begin{equation}
 B_q=\{c\in V:q(c)\ne q(c^*)\},\qquad
 A_e=\{c\in B_q:D(c,e)\ne D(c^*,e)\}.
 \label{eq:kill}
\end{equation}
Truthful outcomes for an experiment set $E$ certify $q(c^*)$ exactly when
\begin{equation}
 \bigcup_{e\in E}A_e=B_q.\label{eq:cover}
\end{equation}
Consequently the minimum target-relative certificate size is the minimum set-cover size; fixed nonnegative experiment costs give weighted set cover. Stored outcomes replace $B_q$ by its uncovered subset before further experiments are chosen.
\end{proposition}
\begin{proof}
The query-disagreeing survivors are exactly $B_q\setminus\bigcup_{e\in E}A_e$. The true model survives, so the query is certified if and only if this set is empty.
\end{proof}

An artifact can store an experimental outcome, a structured model that entails it, or the update rule that obtains its consequence; these representations exchange information according to the distinctions they remove and the operations they preserve. Sound transfer between two representations is the reduced-product construction of abstract interpretation \citep{cousot1979}; \cref{app:reduction} gives its version-space realization.

\begin{theorem}[Adaptive experiments on a pooled memory]\label{thm:policy}
Let a policy choose each experiment and its final answer as functions of a persistent state, and let that state be the pooled outcome $\mu(s)$ of \eqref{eq:independent}. Then the inference state $\delta(\mu(s))$ depends only on the set of outcomes received, and the number of experiments that change it along any branch is at most
\begin{equation}
 \height(\R)-1\le\sum_{j=1}^d\bigl(\height(E_j)-1\bigr).\label{eq:policy}
\end{equation}
On a coupled $W$-coordinate $b$-value memory the count is at most $b^W-1$. Both bounds are attained.
\end{theorem}
\begin{proof}
By \cref{thm:canonical}, $\delta(\mu(s))=[V(s)]$, and $V(s)$ is an intersection over the outcome set. Successive states form a descending chain in $\R$, whose strict steps number at most $\height(\R)-1$; \cref{thm:exchange} bounds this by the channel heights, and \eqref{eq:cardinality} bounds it by $b^W-1$. The chain constructions of \cref{cor:width}, read as policies that eliminate one boundary per experiment, attain both.
\end{proof}

Finite acyclic structural causal models give an explicit intervention instance \citep{pearl2009}. Write $x_i=f_i(x_{\mathrm{pa}(i)},u_i)$ in topological order. Each $f_i$ is a finite feedforward map, and $\doop(x_J=a_J)$ replaces the selected equations by constants. Compiling these maps and replacements gives an attention--ReLU program whose intervention answers agree with the structural equations by induction over the graph. \Cref{prop:cover} then characterizes exactly which stored or newly observed intervention outcomes determine a query in a finite model population. \Cref{app:causal} supplies the construction and a separation between observational agreement and intervention disagreement.

A decoder with unrestricted memory obeys instead $B+t\log_2q\ge N\log_2q$ for tables $f:[N]\to[q]$ reconstructed from a $B$-bit target-dependent artifact and at most $t$ adaptive queries (\cref{app:tables}): stored and queried bits exchange one for one because the transcript is retained. A pooled memory retains only its state, so an experiment counts when it moves the state, and memory rungs and informative rounds are the same quantity. The bit law is the coupled case with unrestricted memory. The observation language and model grammar determine which distinctions require table memory: the AES S-box is specified by finite-field inversion and an affine transformation \citep{nist2023aes}, and that grammar, not the table, is the population.

\section{Computational verification and implications}\label{sec:checks}
Most results are formalized in Lean~4 with Mathlib. The constructions are additionally verified exhaustively on finite families. The network artifacts are compiled from the specified grammars. \Cref{tab:checks} reports the executed checks; \cref{app:checks} defines every enumeration and its assertions.

\begin{table}[t]
\caption{Executed verification. Counts refer to enumerated mathematical objects or input cases, not independent training runs. Every assertion passed.}
\label{tab:checks}
\centering
\small
\begin{tabular}{@{}>{\raggedright\arraybackslash}p{.47\linewidth}r >{\raggedright\arraybackslash}p{.30\linewidth}@{}}
\toprule
Construction & Cases & Verified property\\
\midrule
Meet-closed stores with readouts & 2,120 & Direct and iterative quotients agree\\
Neural transition-table extraction & 72,911 & Every transition recovered exactly\\
Commutative monoids, sizes 1--4 & 106 & Idempotent lifts; 329 chain quotients\\
Chain structural quotients, sizes 1--7 & 127 & Exact residual minimum\\
Compiled attention forests & 356 & All 256 Boolean maps on eight inputs realized\\
Hard code-extraction input cases & 2,848 & Exact semantic round trip\\
Soft code-execution input cases & 8,544 & Maximum absolute error $0$\\
Shared-block transformer executions & 11,392 & Exact outputs and state invariants\\
Symbol-level hardening cases & 1,080 & Exact criterion, including 210 same-symbol tie cases\\
\bottomrule
\end{tabular}
\end{table}

Each transition artifact undergoes a hidden-unit permutation before extraction, and all 11,392 shared-block runs preserve protected coordinates and exact symbolic payloads; \cref{app:checks} defines each enumeration and its assertions.

These results suggest a direct target for learning on artifacts: recover the inference quotient, its decompositions, and the maps between them. The quotient gives a representation-independent object for comparison; the exchange law distinguishes architectures with the same state cardinality; and the hardening certificate identifies transformations that preserve all future inference. In a model population, the extracted object supports questions about which observations separate models, which queries are already determined, and which additional experiments change the answer.

\paragraph{Certified artifact transformations.}
The extracted machine supports three certified operations: state compression, which merges states with the same continuation behavior; representation factorization, which chooses a structural congruence $\alpha$ and computes its minimum compatible residual through \eqref{eq:comp}; and hardening, which preserves the entire transition system through the local symbol-mass certificate. Each is invariant under relabeling of the extracted states, so artifacts are compared by the inference operations they implement and the distinctions those operations retain.

\paragraph{Limitations.}
The bounds measure persistent inference states, not parameter count or running time. Finite-state extraction has exponential worst-case cost in encoded width, and the computational study verifies compiled constructions; their prevalence under optimization is a separate question. Intervention certificates are relative to the declared model population and experiment semantics. The inference theorems require correctness over the specified observation language; obtaining such certificates for learned continuous-state representations is the next empirical and verification problem.

\paragraph{Conclusion.}
An artifact's inference capacity is its retained distinctions together with its composition law: cardinality for a coupled store, chain height for an independent one, and the long revision budget at logarithmic width is the capacity in-context evidence draws on. Attention and feedforward computation realize both laws, finite-symbol hardening preserves their semantic state, and model extraction, representation exchange, and artifact comparison become one operational problem.

\label{mainend}
\clearpage
\bibliographystyle{plainnat}
\bibliography{references}
\clearpage
\appendix
\crefalias{section}{appendix}
\crefname{appendix}{appendix}{appendices}
\Crefname{appendix}{Appendix}{Appendices}
\section{Additional algebraic results}\label{app:algebra}
\subsection{Continuation equivalence and the minimal-state certificate}
We record the details behind the semantic quotient. For every $U,V,H,W\in\Latt$, if $U\equiv_\rho V$, then
\[
 \rho((U\cap W)\cap H)=\rho(U\cap(W\cap H))
 =\rho(V\cap(W\cap H))=\rho((V\cap W)\cap H).
\]
Thus $U\cap W\equiv_\rho V\cap W$. Applying this twice shows that the meet on equivalence classes is independent of both representatives. Associativity, commutativity, idempotence, and the identity class $[\C]$ descend from intersection. The readout is well-defined because the empty continuation is permitted.

If a correct machine reaches the same state after $s,t$, determinism gives equal future outputs after every $u\in P^*$. Writing $H=V(u)$ proves $[V(s)]=[V(t)]$. Conversely, the quotient machine computes the target because its state after $s$ is exactly $[V(s)]$. These observations prove both the lower bound and the existence of a smallest implementation. The argument uses the full declared continuation language $P^*$; the readout need not identify individual candidates.

A finite certificate for a proposed implementation consists of its reachable transition table, a map $\delta$ into $\R$, and the three equations in \eqref{eq:simulation}. Verifying the equations for every reachable state and every input symbol proves correctness for all histories by induction. With $N$ reachable states and $m$ input symbols, the transition part has $Nm$ equalities. This is an all-horizon certificate rather than a test on histories up to a chosen length.

\subsection{Idempotent lifting in finite commutative monoids}\label{app:monoid}
For completeness, let $x$ belong to a finite monoid. There exist positive integers $a,b$ with $x^a=x^{a+b}$. Hence $x^t=x^{t+b}$ for all $t\ge a$. Choose a multiple $r$ of $b$ with $r\ge a$. The exponents $r,2r$ differ by a multiple of $b$, so $x^{2r}=x^r$. Thus $x^r$ is idempotent.

In a commutative monoid, the product of two idempotents is idempotent. On its idempotents, $e\le f$ if $ef=e$ defines a partial order. Reflexivity is idempotence; antisymmetry follows from commutativity; and if $ef=e$, $fg=f$, then $eg=(ef)g=e(fg)=ef=e$. The product $ef$ is the greatest lower bound of $e,f$, and the monoid identity is top.

Let $\pi:S\to R$ be a surjective monoid homomorphism from a finite commutative monoid to a meet-semilattice with top. Since $\pi(x^r)=\pi(x)$ for every $r\ge1$, the idempotents of $S$ still map onto $R$. Given $r_1<\cdots<r_h$, choose idempotent preimages $e_i$ and let $z_i=e_i\cdots e_h$. Their images are $r_i$, and $z_iz_{i+1}=z_i$. This proves the chain-lifting assertion used in \cref{thm:exchange}.

The chain height of a product of finite posets satisfies
\[
 \height\left(\prod_{j=1}^dQ_j\right)
 \le 1+\sum_{j=1}^d(\height(Q_j)-1).
\]
Along a strict chain, coordinate $j$ can change strictly at most $\height(Q_j)-1$ times, and every strict step changes some coordinate. For finite posets the upper bound is attained by interleaving longest chains in the factors. Restricting to a subposet cannot increase height. These facts complete the proof for pooled implementations with arbitrary redundant states.

\subsection{Subdirect representations and residual size}\label{app:congruence}
\begin{proposition}\label{prop:congruence}
Let $a:R\to A$ and $g:R\to G$ be meet homomorphisms on a finite meet-semilattice $R$, with kernels $\alpha,\beta$. The paired representation determines every state of $R$ if and only if $\alpha\cap\beta=\Delta_R$. If $R$ is the minimal inference state, a joint readout correct after every continuation exists if and only if this condition holds.
\end{proposition}
\begin{proof}
Two states have the same pair precisely when they are equivalent in both kernels. Hence the kernel condition is equivalent to injectivity. Injection defines the decoded state and therefore a correct readout. Conversely, if distinct minimal states share a pair, they continue to share it after every meet update because both views are homomorphisms. Minimality supplies a distinguishing continuation, contradicting correctness.
\end{proof}
Every homomorphic image of $R$ is isomorphic to its kernel quotient. Conversely, each congruence quotient supplies a compositional view. The set in \eqref{eq:comp} is nonempty because the equality congruence is feasible. This proves that the minimum is exact, not merely a lower bound. For $\alpha'\subseteq\alpha$, every $\beta$ feasible for $\alpha$ is feasible for $\alpha'$, proving monotonicity under structural refinement.

For a static residual $g:R\to G$ with no compositional requirement, injectivity of $(a,g)$ requires at least $|F|$ residual values on each structural fibre $F$. Let $m$ be the largest fibre size and inject every fibre separately into $[m]$. The resulting residual attains $m$, proving \eqref{eq:static}. It encodes a one-shot reconstruction of the state.

For a chain, every equivalence relation with interval blocks is a meet congruence: if $x\sim y$, taking minimum with $z$ either preserves the pair, maps both to $z$, or produces two points between $x,y$. The interval property gives equivalence in each case. Therefore every choice of adjacent cuts defines a congruence and all chain congruences arise this way. If the structural view cuts a set $A$ of boundaries, the residual cuts must contain its complement. They require $n-|A|=n-a+1$ blocks, and no additional cuts are necessary.

The Boolean coordinate law follows without requiring a normalized residual. Fix a structural fibre and choose a reaching history for each of its $2^{k-s}$ states. Correctness and \cref{thm:canonical} require distinct joint implementation states. Since their structural values agree, the residual values are pairwise distinct. This lower bound applies to any finite independently pooled residual. The complementary-coordinate projection attains it.

\subsection{A presentation with many survivors but one inference state}
Let $\C=\{c_0,c_1,\ldots,c_k\}$, let observations independently remove any of $c_1,\ldots,c_k$, and let no observation remove $c_0$. The reachable version spaces are all subsets containing $c_0$, so $|\Latt|=2^k$. First-fit selection in the displayed order always returns $c_0$, including after every continuation. Hence $|\R|=1$. This example makes the readout dependence of the lower bound explicit. Counting reachable survivor sets is sharp precisely when the continuation equivalence separates them.

\section{Extraction and finite abstractions}\label{app:extraction}
\subsection{Operational reconstruction of the meet}
\begin{proposition}[Recovering the inference algebra]\label{prop:recovermeet}
Let $(Q,q_0,(t_p)_{p\in P},o)$ be a finite reachable deterministic Moore machine. Suppose the transition maps commute and are idempotent:
\[
 t_pt_q=t_qt_p,\qquad t_p^2=t_p\quad(p,q\in P).
\]
The operation $r\wedge s=t_{w_s}(r)$, for an access word $w_s$ of $s$, is a meet with top $q_0$ satisfying
\[
 t_p(q)=q\wedge t_p(q_0),
\]
and any meet operation with top $q_0$ satisfying this identity equals it. Uniqueness uses reachability and the identity alone; commutativity and idempotence enter through existence. The operation can be reconstructed from the transition table. Its version-space realization on candidate set $Q$ has consistency sets $K(p)=\mathord{\downarrow}t_p(q_0)$ and survivor state $\mathord{\downarrow}q_s$ after history $s$.
\end{proposition}
\begin{proof}
For a word $w$, write $t_w$ for its transition map. Commutativity and idempotence of the generators imply these properties for all $t_w$. For each state $r$, choose an access word $w_r$ with $t_{w_r}(q_0)=r$. Define $r\wedge s=t_{w_s}(r)$. If $w'_s$ is another access word for $s$, then, using an access word for $r$,
\[
 t_{w_s}(r)=t_{w_s}t_{w_r}(q_0)
 =t_{w_r}(s)=t_{w_r}t_{w'_s}(q_0)=t_{w'_s}(r).
\]
Thus the definition is independent of the access word. It is commutative by the same calculation. The concatenation $w_sw_t$ is an access word for $s\wedge t$, which proves associativity. Idempotence follows from $t_{w_r}^2=t_{w_r}$, and the empty word gives identity $q_0$. Therefore the operation is a meet with top.

Taking the one-letter access word for $t_p(q_0)$ gives the transition identity. In any other meet representation with that identity, induction on an access word forces $r\wedge s=t_{w_s}(r)$, using only reachability and the identity; this proves uniqueness. Finally,
\[
 \mathord{\downarrow}r\cap\mathord{\downarrow}s
 =\mathord{\downarrow}(r\wedge s),\qquad
 \mathord{\downarrow}q_0=Q.
\]
Induction on observations gives the stated survivor sets. The map $r\mapsto\mathord{\downarrow}r$ is injective, and the readout is $\rho(\mathord{\downarrow}r)=o(r)$.
\end{proof}

For a minimal machine whose output on histories is invariant to permutation and repetition, the transition maps necessarily commute and are idempotent. Indeed, applying $pq$ or $qp$ after any reaching history, followed by any continuation, gives equal outputs; minimality makes the resulting states equal. The same argument compares $pp$ with $p$. Conversely, commuting idempotent transitions imply both invariances directly. Consequently the condition in \cref{prop:recovermeet} can be tested from the extracted minimal machine, without presupposing a model population.

The resulting model is operational. It identifies the distinctions implemented by the artifact under its observation interface. Connecting those states to external hypotheses, causal mechanisms, or object identities requires a semantic interpretation of the observations and readout. When such an interpretation is supplied by \eqref{eq:version}, \cref{thm:canonical} makes the correspondence unique on reachable states.

\subsection{Partition refinement and certificates}
Let $N$ be the number of reachable states. Start with equivalence of output labels and refine by
\[
 q\sim_{r+1}q'
 \quad\Longleftrightarrow\quad
 o(q)=o(q')\text{ and }t_p(q)\sim_rt_p(q')\text{ for all }p\in P.
\]
Induction shows that $\sim_r$ is agreement on all continuations of length at most $r$. Each strict refinement increases the number of blocks, so stabilization occurs after at most $N-1$ strict refinements. At stabilization the relation is a congruence for every transition. Induction on arbitrary continuations then proves agreement at every length. For two different blocks, the refinement history yields a finite distinguishing word. Thus a complete extraction certificate consists of the transition table, the quotient projection, and distinguishing continuations between distinct blocks.

For a finite input alphabet, testing commutativity and idempotence requires finitely many state equalities. Once \cref{prop:recovermeet} applies, its access-word construction gives the meet table. Enumerating finite set partitions and testing meet compatibility enumerates every meet congruence. Evaluating \eqref{eq:comp} is therefore an exact finite procedure on an extracted artifact. These quantities are invariant under isomorphisms of the extracted machine.

\subsection{Continuous-state artifacts with finite invariant regions}
Exact extraction is not confined to one-hot activations. Let $Z\subseteq\RR^d$ be an invariant state set and let $B_1,\ldots,B_N$ be a finite partition of $Z$. Assume that each input $p$ maps every region into one region and that the readout is constant on each region:
\begin{equation}
 F_\theta(B_i,p)\subseteq B_{T(i,p)},\qquad
 o_\theta(z)=o_i\quad(z\in B_i),\qquad z_0\in B_{i_0}.
 \label{eq:regions}
\end{equation}
\begin{proposition}[Finite-region extraction]\label{prop:regions}
The finite machine $(T,i_0,o)$ has exactly the same output on every history as the continuous-state artifact. If $Z$ and its state regions have rational semialgebraic descriptions, the initial state and input encodings are rational, output labels have rational encodings, and the transition and readout consist of rational affine, ReLU, and deterministically tie-broken hard-attention operations, all conditions in \eqref{eq:regions}, including coverage and disjointness, are decidable.
\end{proposition}
\begin{proof}
The initial state belongs to the declared initial region. If the current state belongs to $B_i$, the inclusion in \eqref{eq:regions} puts its successor in $B_{T(i,p)}$; the readout equality then proves the simulation. Induction gives equality for every history.

Affine and ReLU operations have semialgebraic graphs. A selected argmax index $j$ is described by the finite conjunction $s_j>s_i$ for $i<j$ and $s_j\ge s_i$ for $i>j$. Hence a fixed hard-attention network has a semialgebraic graph. Each inclusion or readout equality in \eqref{eq:regions} is a first-order sentence over the ordered real field with rational coefficients. Real quantifier elimination decides it \citep{basu2006}. Coverage and disjointness are sentences of the same form.
\end{proof}
The proposition gives a certificate checker for a supplied abstraction. The same partition-refinement procedure extracts its minimal machine. Region labels need not be the final semantic states: minimization removes redundant distinctions between regions.

\section{A literal attention--ReLU compiler}\label{app:compiler}
\subsection{Leaf maps and transition maps}
For a map $f:[m]\to[q]$, set the $i$th column of $H_f$ to $e_{f(i)}$. Then $H_fe_i=e_{f(i)}$. A ReLU layer with identity input weights is the identity on nonnegative one-hot inputs, so a standard two-affine-layer feedforward block implements the table exactly. For a map on two finite inputs, the units $\relu(z_i+p_j-1)$ in \eqref{eq:ffntable} are zero except at the unique active pair. Their output matrix therefore gives the exact transition table. The construction uses $Nm$ hidden units and rational weights in $\{-1,0,1\}$ for its conjunction and one-hot output stages.

A leaf table is an operational parameter matrix. The extractor obtains it by evaluating the leaf on every basis input, whether or not the implementation's hidden units have been permuted or reorganized. The same procedure extracts a transition table from any implementation satisfying the finite interface, not just from the displayed compiler.

\subsection{Tuple formation and finite dataflow}
To combine $r$ already computed symbols, use one fixed-mask attention head per input to retrieve its one-hot vector $v^{(1)},\ldots,v^{(r)}$ into disjoint scratch blocks. For a tuple $(i_1,\ldots,i_r)$, the hidden unit
\begin{equation}
 h_{i_1,\ldots,i_r}=\relu\!\left(\sum_{a=1}^r v^{(a)}_{i_a}-(r-1)\right)
 \label{eq:tuple}
\end{equation}
is one exactly at the active tuple and zero elsewhere. An output matrix therefore implements any finite map of the retrieved tuple. Shared node-specific dispatch uses the unit
$\relu(u_v+\sum_a v^{(a)}_{i_a}-r)$ instead. A topological schedule of these blocks implements every finite acyclic dataflow program. Fixed-mask reads have a single allowed key, so their outputs are identical for hard and soft attention. In particular, the tuple construction composes the forest outputs and implements parent retrieval and structural equations in \cref{app:causal}; no external tuple-encoding oracle is needed.

\subsection{Hard attention gates}
Write each routing score as
\[
 s_j(x)=\alpha_j^\top\phi(x)+\beta_j
 =\langle (\phi(x),1),(\alpha_j,\beta_j)\rangle.
\]
The child tokens carry key $(\alpha_j,\beta_j)$ and value $e_{\dec{c_j}(x)}$. The parent query is $(\phi(x),1)$. The usual attention scale can be absorbed into the keys. With the declared deterministic tie rule, hard attention returns the value of $j^*(x)$. Multiplication by $H_g$ gives the node's value. Induction from leaves proves the semantic identity at every node.

The extractor reads the effective query--key score coefficients, the wiring or mask, and the finite table induced by the feedforward block. Different factorizations into query and key matrices can give the same scores; the extracted code is determined semantically rather than by selecting one parameter factorization. Child reorderings preserve the extracted semantics when the tie convention is transported with the ordering. The equality in \eqref{eq:roundtrip} concerns the decoded code, not syntactic equality with its source.

\subsection{Shared positionwise feedforward maps and residual paths}
We give a standard masked-transformer realization of the forest dataflow. Unfold a finite directed acyclic graph into a finite forest when a node has multiple parents. Use one token for each resulting node. Embed the finite node-output alphabets into a common tagged alphabet, so every payload uses a common coordinate block. The token coordinates contain: a fixed one-hot node identifier; protected input features $(\phi(x),1)$; the node's key for its parent; a one-hot payload; and a scratch block initialized to zero. The identifiers and keys are operational inputs to shared dispatch and routing.

Process nodes by height above the leaves. In a layer, the mask of each active parent exposes exactly its children. Inactive tokens attend to themselves. Shared query and key projections select the protected query and key blocks. The shared value projection reads the payload; the output projection writes into scratch. Because scratch was zero, the attention residual leaves there exactly the attended value. Protected coordinates and the old payload remain unchanged.

For hard attention, scratch is one-hot. Let $u_v$ be the identifier of node $v$, and let $w_j$ be the scratch symbol coordinate. The shared hidden unit
\[
 h_{vj}=\relu(u_v+w_j-1)
\]
is active precisely for token $v$ and its selected symbol $j$. An output matrix maps it to the payload required by that node's finite map. Inactive nodes instead retain their old payload using the same conjunction construction with old-payload coordinates. The feedforward residual update subtracts the old payload and adds the desired one, and subtracts scratch to reset it to zero. Subtraction of a nonnegative coordinate is implemented by a negative output weight on its ReLU. All protected coordinates receive zero update. Thus one shared positionwise feedforward map implements every node-specific operation at that layer.

A preliminary finite-map layer initializes leaf payloads from the one-hot input and initializes internal payloads to a fixed dummy symbol. Hence every payload is a valid one-hot vector from the start. Induction over heights proves that active children have already acquired the correct symbol and that completed nodes retain it. A forest of depth $D$ needs at most $D$ routing layers after this initialization, followed by the declared root readout. This construction fixes all masks, projections, residual paths, and feedforward dispatch; it introduces no unused field encoding the original program identifier.

\subsection{Finite-symbol correction within the shared block}
Correction can also be implemented in a shared single-hidden-layer feedforward block. For a token identifier $u_v\in\{0,1\}$ and scratch coordinate $w_j\in[0,1]$, define
\begin{equation}
 h_{vj}^{\eta}=
 \frac{\relu(w_j+u_v-1-\eta)-\relu(w_j+u_v-2+\eta)}{1-2\eta}.
 \label{eq:gatedround}
\end{equation}
If $u_v=1$, this equals $R_\eta(w_j)$. If $u_v=0$, both ReLUs vanish because $w_j\le1$. The output matrix can therefore apply the node-specific finite map to the corrected symbol. The payload and scratch overwrite construction is unchanged. Whenever \eqref{eq:exactmass} holds, the outgoing payload is exactly one-hot at each layer. This gives a literal implementation of \cref{thm:hardening}, with shared parameters across token positions.

For real inputs beyond a finite input alphabet, the routing induction remains valid on any domain on which leaf maps have the stated implementation and the gate conditions hold. In the finite construction, all feature values and leaf maps are part of the specified input embedding, so every operation used in extraction is explicitly available.

\section{Hardening, perturbations, and observability}\label{app:hardening}
\subsection{Sharpness of the symbol-mass bound}
Let the total unnormalized weight of the target symbol be $G$ and that of all other symbols be $B$. Then $\eps_{\ne y^*}=B/(G+B)$. Under the hypotheses of \cref{thm:hardening}, $G\ge m_*e^{s_*/T}$ and $B\le k_-e^{(s_*-\delta)/T}$, yielding \eqref{eq:margin}. Equality is attained when the $m_*$ good branches all have score $s_*$, all $k_-$ bad branches have score $s_*-\delta$, and no other branches are present. Solving $r/(1+r)\le\eta$ gives $r\le\eta/(1-\eta)$ and the unique-winner temperature bound \eqref{eq:temperature}.

The local criterion in \eqref{eq:exactmass} is necessary and sufficient for correction to the selected one-hot symbol. A later finite map can collapse several symbols and thereby permit further equivalences; the criterion is applied before that map. It already accounts for branch equivalence at the symbol level. In particular, if every branch returns the same symbol, $\eps_{\ne y^*}=0$ for all temperatures, regardless of routing ties.

\subsection{A perturbation certificate}
\begin{corollary}[Numerical robustness of correction]\label{cor:robust}
Let the ideal attention mixture have wrong-symbol mass at most $\eta-\xi$ for some $\xi>0$. If its numerical realization differs coordinatewise by at most $\xi$, then applying $R_\eta$ still returns the exact selected one-hot symbol.
\end{corollary}
\begin{proof}
The selected coordinate is at least $1-\eta+\xi$ before perturbation and at least $1-\eta$ afterward. Every other coordinate is at most $\eta-\xi$ before perturbation and at most $\eta$ afterward. The correction function is constant on these two regions.
\end{proof}
If each score is perturbed by at most $\zeta$ and a unique winning score has gap $\delta>2\zeta$, the winner is unchanged and the perturbed gap is at least $\delta-2\zeta$. Substituting this gap into \eqref{eq:margin}, followed by \cref{cor:robust}, gives a score-and-value perturbation certificate. The same-symbol version permits route changes within a group whose values coincide; it suffices to preserve a positive gap from that group's largest score to every different-symbol branch.

\subsection{Approximate hardening without symbolic correction}
For bounded expert values $\norm{v_j}\le B$ and a unique winning score with gap $\delta$, put $r=(k-1)e^{-\delta/T}$. Directly,
\begin{equation}
 \left\|\sum_jp_j(T)v_j-v_{j^*}\right\|
 \le 2B\frac{r}{1+r}\le 2B(k-1)e^{-\delta/T}.
 \label{eq:approxhard}
\end{equation}
For compositions $F_L^T\circ\cdots\circ F_1^T$ and $F_L^0\circ\cdots\circ F_1^0$, assume that a common region contains all compared trajectories, each $F_\ell^T$ is $L_\ell$-Lipschitz there, and $\sup_z\norm{F_\ell^T(z)-F_\ell^0(z)}\le\eps_\ell(T)$. Inserting and subtracting $F_\ell^T(z_{\ell-1}^0)$ yields
\[
 e_\ell\le L_\ell e_{\ell-1}+\eps_\ell(T),\qquad
 e_L\le\sum_{\ell=1}^L\eps_\ell(T)\prod_{j=\ell+1}^LL_j.
\]
Uniform score separation is imposed on the region where the local hard/soft bound is used. Alternatively, a stepwise perturbation argument certifies preservation of each gate's winning symbol. These hypotheses expose the dependence of approximate hardening on downstream stability; exact symbolic correction removes that accumulation by restoring a codebook element after each gate.

Let $\widehat L_T(\theta)$ and $\widehat L_0(\theta)$ be empirical losses on a fixed sample. If $|\widehat L_T(\theta)-\widehat L_0(\theta)|\le\eps$ for every $\theta$ in a declared parameter set, then their approximate-fit sets obey
\[
 \{\theta:\widehat L_T(\theta)\le a-\eps\}
 \subseteq\{\theta:\widehat L_0(\theta)\le a\}
 \subseteq\{\theta:\widehat L_T(\theta)\le a+\eps\}.
\]
This follows by adding the uniform loss discrepancy. Under exact symbolic preservation and a loss on the resulting symbols, the two sets are equal for every threshold.

\subsection{Artifact observations and their sufficiency}
For maps $O:\Theta\to\mathcal O$, $A:\Theta\to\mathcal A$, and $Q:\Theta\to\mathcal Q$, an exact inference function on the image of $(O,A)$ exists if and only if $Q$ is constant on its fibres. Necessity follows by applying the inference function to equal observation pairs. For sufficiency, assign to each pair in the image the common value of $Q$ on its fibre. This defines a function with the required factorization and proves \eqref{eq:sufficiency}.

If $O_0=h\circ O_T$, equality of $O_T$ implies equality of $O_0$, proving \eqref{eq:coarsen}. This is an observational comparison, with a declared coarsening map. It does not equate the activations of independently run soft and hard networks when their upstream states change. \Cref{thm:hardening} supplies the separate execution-equivalence theorem for the compiled symbolic interfaces.

\section{Sound exchange between representations}\label{app:reduction}
Let $\mathcal A,\mathcal G\subseteq\powerset(\C)$ be families containing $\C$ and closed under arbitrary intersections. Since $\C$ is finite, it suffices to require closure under finite intersections. Define
\[
 \cl_{\mathcal A}(V)=\bigcap\{A\in\mathcal A:V\subseteq A\},\qquad
 \cl_{\mathcal G}(V)=\bigcap\{G\in\mathcal G:V\subseteq G\}.
\]
A structural state $A$ and residual state $G$ jointly denote $A\cap G$. Their reduced product \citep{cousot1979} is
\[
 \operatorname{red}(A,G)=
 \bigl(\cl_{\mathcal A}(A\cap G),\cl_{\mathcal G}(A\cap G)\bigr).
\]
\begin{proposition}
Reduction is componentwise contracting, preserves the joint version space exactly, and is idempotent.
\end{proposition}
\begin{proof}
Write $V=A\cap G$, $A'=\cl_{\mathcal A}(V)$, and $G'=\cl_{\mathcal G}(V)$. Since $A$ and $G$ are available overapproximations, $V\subseteq A'\subseteq A$ and $V\subseteq G'\subseteq G$. Hence $V\subseteq A'\cap G'\subseteq A\cap G=V$. Repeating reduction uses the same $V$ and therefore produces the same pair.
\end{proof}
The representation in one component can sharpen the other by making consequences of their joint information expressible there. No additional external observation is needed for this transfer. The joint set of possible models is unchanged, so the operation preserves evidence.

\section{Model queries and causal interpretation}\label{app:causal}
\subsection{Query certificates and adaptive experiments}
Fix $c^*\in V$. After truthful outcomes on $E$, the survivor set is
\[
 V_E=\{c\in V:D(c,e)=D(c^*,e)\text{ for all }e\in E\}.
\]
The true candidate belongs to this set. Intersecting with $B_q$ gives $B_q\setminus\bigcup_{e\in E}A_e$, proving \cref{prop:cover}. For a stored experiment set $E_0$, an additional set $E_1$ certifies the query exactly when its kill sets cover $B_q\setminus\bigcup_{e\in E_0}A_e$. Model identification is the special case $q(c)=c$, or a semantic equivalence class of $c$ when extensional duplicates are identified.

The result is target-relative. A policy choosing experiments before learning which model is true is represented by a decision tree whose branches are outcomes. Every root-to-leaf path must satisfy the same covering condition for each true model reaching that leaf. Thus the proposition characterizes the certificates required at policy leaves; optimizing the entire adaptive policy additionally chooses how those certificates share prefixes.

\subsection{Compiling a finite acyclic structural causal model}
Let the endogenous variables $X_1,\ldots,X_d$ and exogenous variable $U$ have finite alphabets. For each model $c$, let its graph be acyclic, with structural maps
\[
 X_i=f_{c,i}(X_{\mathrm{pa}_c(i)},U).
\]
For a fixed intervention $I=\doop(X_J=a_J)$ and exogenous input $u$, evaluate the variables in a topological order, using the constant $a_i$ for $i\in J$ and the structural table otherwise. Each operation is a finite map. Retrieve its already evaluated parent symbols and the exogenous symbol with fixed-mask attention heads. The conjunction units in \eqref{eq:tuple} form the tuple basis vector, and the table matrix applies its structural map. An intervention-controlled finite map selects the constant branch or the structural branch. The forest compiler realizes these operations by attention and ReLU maps.

\begin{proposition}[Intervention preservation]\label{prop:scm}
For every model $c$, intervention $I$, and exogenous input $u$, the compiled program returns the unique solution of the intervened structural equations. For a specified finite exogenous distribution, weighted aggregation gives the exact intervention distribution.
\end{proposition}
\begin{proof}
At the first variable in a topological order, the program applies either its intervention constant or the structural map with exogenous input $u$. This agrees with the intervened equation. Suppose the program is correct for all preceding variables. Every parent of the next variable precedes it, so the retrieved tuple is correct; applying its selected table proves the next value correct. Induction gives the unique solution. Summing its one-hot output distribution against the specified exogenous probabilities gives the exact interventional law.
\end{proof}
To apply the deterministic certificate theorem, an experiment can report a controlled-exogenous outcome or the exact observable distribution under an intervention. The output of such an experiment is a deterministic function of the candidate model. Observations from a stochastic experiment require their own statistical consistency map; the survivor and continuation constructions then use that declared map.

For an explicit separation, let $U$ be uniform on $\{0,1\}$ and consider
\[
 c_1:\ X=U,\ Y=X,\qquad c_2:\ X=U,\ Y=U.
\]
Both give the same observational distribution: probability $1/2$ each on $(0,0)$ and $(1,1)$. Under $\doop(X=0)$, the first model has $Y=0$ surely, while the second retains uniform $Y$. The query $\Pr(Y=1\mid\doop(X=0))$ is therefore $0$ or $1/2$. Observational distribution access does not separate the pair, while that intervention distribution does. Their structural tables, or a compiled program with the same intervention semantics, also separate the query by \eqref{eq:sufficiency}.

\section{Population memory bounds and routed model classes}\label{app:tables}
\begin{proposition}[Table reconstruction]
Suppose a fixed deterministic decoder reconstructs every function $f:[N]\to[q]$ from a target-dependent $B$-bit artifact and at most $t$ adaptive table queries. Then
\[
 B+t\log_2q\ge N\log_2q.
\]
For a population of all permutations of $[N]$, with $0\le t\le N$, the corresponding bound is
\[
 B\ge\log_2((N-t)!).
\]
\end{proposition}
\begin{proof}
Fix the artifact. Its adaptive query tree has at most $q^t$ leaves after padding shorter transcripts to depth $t$. The artifact and answer transcript determine all queries and the reconstructed table. Two different tables cannot have the same pair, since the decoder is exact. There are at most $2^Bq^t$ pairs and $q^N$ tables, giving the first inequality.

For permutations, repeated queries supply no new information and can be omitted. Pad to $t$ distinct queries. The possible answer counts at successive levels are at most $N,N-1,\ldots,N-t+1$, independent of how query locations are chosen. There are at most $N!/(N-t)!$ leaves for each artifact. Comparing the $2^B N!/(N-t)!$ pairs with the $N!$ targets proves the second inequality.
\end{proof}
All target-dependent information, including a varying architecture, routing parameters, learned encoders, and finite tables, belongs in the $B$-bit description. The bound concerns a population with a fixed decoder, not the shortest program of a particular function. A restricted algebraic population can have a shorter description. This is why a structured S-box implementation and an arbitrary permutation table belong to different counting problems.

\subsection{Selection and execution have different statistical complexity}
Let $\mathcal F$ consist of routed functions $f(x)=h_{r(x)}(x)$, with routing class $\mathcal H_{\mathrm{route}}$ and expert classes $\mathcal H_i$, all with finite output alphabets. For a finite sample $S$, let $S_i(r)$ be the sample positions routed to expert $i$. Then
\begin{equation}
 \bigl|\mathcal F\!\restriction_S\bigr|
 \le\sum_{r\in\mathcal H_{\mathrm{route}}\restriction_S}
          \prod_i\bigl|\mathcal H_i\!\restriction_{S_i(r)}\bigr|.
 \label{eq:localgrowth}
\end{equation}
For each realized routing pattern, the output restriction is determined by the expert restrictions on their assigned positions. Multiplying counts bounds that pattern; summing over patterns gives \eqref{eq:localgrowth}. Fixed experts contribute factors of one. When expert contents also vary, their restrictions contribute explicitly. Thus an economical selection representation does not erase the variability of the models it executes.

\section{Hardening and measurability}\label{app:measurability}
Discontinuity of argmax is not, by itself, a failure of measurable inference. A deterministic finite argmax of measurable scores is measurable: each selected-index event is a finite conjunction of measurable comparisons, with strict comparisons for earlier indices to implement the tie rule. Finite compositions with measurable feedforward maps remain measurable.

\begin{proposition}[Finite parameter families]
For a finite parameter set $\Theta$ and measurable losses $\ell_\theta(z)$ with defined finite expectations $L(\theta)$, the uniform-deviation event
\[
 \left\{(z_1,\ldots,z_m):
 \max_{\theta\in\Theta}\left|L(\theta)-\frac1m\sum_{i=1}^m\ell_\theta(z_i)\right|>\eps\right\}
\]
is measurable. This holds for both soft and deterministically tie-broken hard attention whenever their individual losses are measurable.
\end{proposition}
\begin{proof}
Each empirical loss is measurable in the sample. The maximum is finite, so its threshold event is a finite union of measurable events.
\end{proof}

A second statement applies to real-valued parameters. Fix a finite architecture of affine, ReLU, and hard-attention operations, a semialgebraic parameter set $\Theta\subseteq\RR^p$, and a semialgebraic loss. If the reference distribution has finite support, then the population loss is semialgebraic in the parameters. The graph of the network is semialgebraic by the same argmax description used in \cref{prop:regions}. The joint strict-deviation set in parameter and sample variables is therefore semialgebraic. Its projection to sample variables is semialgebraic by the Tarski--Seidenberg theorem, hence Borel \citep{basu2006}.

For general parameter spaces and joint families, projection raises a separate descriptive-set-theoretic issue: projections of Borel sets need not be Borel \citep{kechris1995}. The relevant object is the parameter-uniform event, not an individual network evaluation. The preservation and observation-coarsening results in \cref{sec:hardening} identify how hardening changes computation and identifiability without identifying either change with a loss of measurability.

\section{Reproducibility and executed checks}\label{app:checks}
All results in this paper are verified in Lean~4 with Mathlib. The accompanying implementation consists of four Python modules and a machine-readable result file. Running
\begin{verbatim}
python checks/verify.py
\end{verbatim}
regenerates every reported count. The recorded environment is Python 3.13.5 with NumPy 2.3.5. Integer checks use exact integer arithmetic. Float64 is used for dot products and softmax evaluation in the finite attention constructions. No fitted parameters, external datasets, or randomized train/test splits enter the enumerations.

\paragraph{Inference quotients.}
For candidate sizes $k=1,2,3$, enumerate every family of bit masks containing the full mask and closed under bitwise intersection. There are 70 such labeled families in total. For each family, enumerate every Boolean readout and add the first-surviving-index readout, with a separate inconsistency output. This gives 2,120 cases. Direct equivalence by all meets is compared against iterative transition refinement; readout descent and meet compatibility are also checked.

\paragraph{Congruences and exact residual minima.}
Enumerate every set partition of each chain of size one through seven, retaining exactly those compatible with minimum. There are $2^{n-1}$ at size $n$ and 127 structural quotients across the range. For each, search all residual congruences for the smallest jointly injective representation. The result is $n-a+1$ in every case. Repeat for Boolean stores on one through three coordinates and every initial coordinate projection; the minima are $2^{k-s}$. The four-chain parity residual is explicitly checked to be injective jointly with the structural view but incompatible with minimum.

\paragraph{Non-idempotent pooling.}
Enumerate every commutative multiplication table with identity labeled zero on carriers of sizes one through four, and retain the associative tables. The counts are $1,2,9,94$, totaling 106 labeled monoids; 93 are non-idempotent. For every element, a factorial power is checked to be idempotent. Enumerate every surjective homomorphism from each monoid to a chain of size no larger than its carrier. All 329 maps are checked to remain surjective on idempotents and satisfy the height bound. These are labeled tables, not isomorphism classes.

\paragraph{Neural machine extraction.}
Compile each of the 2,120 inference machines by \eqref{eq:ffntable}. Apply a hidden-unit permutation using seed 20260907, transporting both adjacent weight matrices. Extract the transition and readout tables from the resulting network and minimize them. All 72,911 state--input evaluations reproduce the specified transition, and every extracted quotient equals the direct continuation equivalence.

\paragraph{Attention forests.}
The input set is $\{0,\ldots,7\}$. The four base maps are constant zero, constant one, parity, and complemented parity. Enumerate all one-gate combinations of these maps over thresholds $1.5,3.5,5.5$, with either the identity or complemented output map. This gives 96 one-gate codes, in addition to four leaves. Finally enumerate all $4^4=256$ base-map assignments to the four leaves of a balanced depth-two tree whose thresholds are $1.5,3.5,5.5$. The total is 356 network artifacts, realizing all 256 Boolean functions on eight inputs.

Routing scores at threshold $t$ are $t-x$ and $x-t$. Their minimum gap on the integer input set is one. The feedforward correction uses $\eta=1/4$ and then applies the chosen finite map. Extracted routing coefficients, leaf tables, and output maps reproduce the code on all 2,848 network--input pairs. Softmax execution at temperatures $0.1$, $0.25$, and $0.5$ gives 8,544 further pairs. The maximum recorded absolute error from the target one-hot output is zero in float64. These equalities check the compiled finite-symbol construction; the mathematical guarantee is \cref{thm:hardening}.

\paragraph{Literal shared-block execution.}
Compile the same 356 forests into token matrices with protected node identifiers, input features, parent keys, payloads, and scratch blocks. Execute the actual shared query/key/value/output projections, masked attention, and shared ReLU feedforward matrices with both residual paths. For each artifact and each of the eight inputs, run hard attention and temperatures $0.1$, $0.25$, and $0.5$. All 11,392 runs return the target symbol with zero maximum recorded absolute error. After every layer, protected coordinates are unchanged, payloads are one-hot, and scratch is zero; the largest recorded error in these invariants is also zero. The forward evaluator receives only matrices, masks, and the input, not the source program.

\paragraph{Same-symbol ties and coupled minimum.}
Enumerate every three-branch score vector in $\{-1,0,1\}^3$, every binary assignment of branch values, and temperatures $0.1,0.25,0.5,1,2$. The 1,080 cases verify the necessary-and-sufficient wrong-symbol-mass criterion; 210 involve tied maximum-score branches with identical symbols. Separately, allocate chain cuts for sizes two through 32 and alphabet sizes two through eight, giving 217 sharp independent-width constructions. The coupled digit-code attention minimum is verified for chain sizes two through 64 and alphabets $2,3,4,8$, covering 357,756 input pairs.

\end{document}
